\documentclass{beamer}

\usetheme{Madrid}
\usecolortheme{default}

\title{Multimodaler Sprach-Assistent}
\subtitle{KI-basierte Sprach-zu-Bild Pipeline}
\author{Ruba Shehadeh, Godfrey Uchechukwu Atulegwu, Kilian Ketelhohn}
\date{\today}

\begin{document}

% Titelfolie
\frame{\titlepage}

% ----------------------------
\begin{frame}{Projektidee}
\begin{itemize}
    \item Entwicklung eines intelligenten Sprach-Assistenten
    \item Verarbeitung gesprochener Sprache in Echtzeit
    \item Transformation von Sprache zu strukturierter Information
    \item Visualisierung durch generierte oder gesuchte Bilder
\end{itemize}
\end{frame}

% ----------------------------
\begin{frame}{Ziel des Projekts}
\begin{itemize}
    \item Verbindung von:
    \begin{itemize}
        \item Spracherkennung (Speech-to-Text)
        \item Natural Language Processing (NLP)
        \item Informationsextraktion
        \item Bildgenerierung / Bildsuche
    \end{itemize}
    \item Interaktive KI-Pipeline für Benutzer
\end{itemize}
\end{frame}

% ----------------------------
\begin{frame}{Systemarchitektur}
\begin{itemize}
    \item Mikrofon-Eingabe
    \item Speech-to-Text (Whisper)
    \item NLP Analyse (spaCy)
    \item Extraktion von Verben und Nomen
    \item Bildgenerierung (Unsplash API)
    \item Darstellung im Streamlit Chat Interface
\end{itemize}
\end{frame}

% ----------------------------
\begin{frame}{Speech-to-Text}
\begin{itemize}
    \item Nutzung von OpenAI Whisper
    \item Wandelt Audio in Text um
    \item Unterstützt verschiedene Audioformate
    \item Grundlage für weitere Verarbeitung
\end{itemize}
\end{frame}

% ----------------------------
\begin{frame}{Natural Language Processing}
\begin{itemize}
    \item Tokenisierung des Textes
    \item Part-of-Speech Tagging
    \item Identifikation von:
    \begin{itemize}
        \item Verben (Handlungen)
        \item Nomen (Objekte)
    \end{itemize}
    \item Verwendung von spaCy
\end{itemize}
\end{frame}

% ----------------------------
\begin{frame}{Informationsextraktion}
\begin{itemize}
    \item Herausfiltern relevanter Begriffe
    \item Fokus auf Nomen als Schlüsselbegriffe
    \item Nutzung für Bildsuche oder KI-Generierung
\end{itemize}
\end{frame}

% ----------------------------
\begin{frame}{Bildgenerierung}
\begin{itemize}
    \item Nutzung von Unsplash API
    \item Suchanfrage basiert auf extrahierten Nomen
    \item Automatische Visualisierung des Textinhalts
\end{itemize}
\end{frame}

% ----------------------------
\begin{frame}{Frontend (Streamlit)}
\begin{itemize}
    \item Interaktive Web-Oberfläche
    \item Upload von Audio oder Texteingabe
    \item Darstellung von:
    \begin{itemize}
        \item Analyseergebnissen
        \item Markierten Verben
        \item Generierten Bildern
    \end{itemize}
\end{itemize}
\end{frame}

% ----------------------------
\begin{frame}{Technologien}
\begin{itemize}
    \item Python
    \item Whisper (Speech-to-Text)
    \item spaCy (NLP)
    \item Streamlit (UI)
    \item Unsplash API (Bildsuche)
\end{itemize}
\end{frame}

% ----------------------------
\begin{frame}{Ergebnis}
\begin{itemize}
    \item Sprachinput wird in strukturierte Daten umgewandelt
    \item Linguistische Analyse in Echtzeit
    \item Visuelle Darstellung der Inhalte
    \item Kombination mehrerer KI-Technologien
\end{itemize}
\end{frame}

% ----------------------------
\begin{frame}{Fazit}
\begin{itemize}
    \item Erfolgreiche Umsetzung einer multimodalen KI-Pipeline
    \item Verbindung von Sprache, Text und Bild
    \item Erweiterbar für zukünftige KI-Anwendungen
\end{itemize}
\end{frame}

% ----------------------------
\begin{frame}{Ende}
\centering
\Huge Fragen?
\end{frame}

\end{document}